%! Setting Up a Test for a Population Mean — Unit 7, Lesson 7.4 — Homework
\documentclass[10pt]{article}
\usepackage{apstats-article}
\usepackage{apstats-boxes}
\newcommand{\TallMath}[1]{$\displaystyle #1\rule[-1.4em]{0pt}{3.2em}$}

\begin{document}

\pageheader{Unit 7, Lesson 7.4}{Homework}
\namedateperiod

\vspace{0.15in}

\textit{For each problem: (a) name the inference method and df; (b) state $H_0$ and $H_a$;
(c) verify both conditions. Show all work.}

\begin{enumerate}

\item \textbf{Water Temperature.} \quad
A marine biologist believes the mean water temperature at a reef has risen above the historical
average of 77$^\circ$F. She randomly collects 28 temperature readings; the data are slightly
right-skewed with no outliers. The sample standard deviation is 2.3$^\circ$F.

\begin{enumerate}[label=\alph*.]
  \item Method: \blank{5.5cm} \quad $df = $ \blank{1.5cm}
  \item $H_0$: \blank{4.5cm} \quad $H_a$: \blank{4.5cm}
  \item Conditions:
        \begin{itemize}
          \item Random: \writeline
          \item Normal ($n < 30$): \writeline
        \end{itemize}
\end{enumerate}
\vspace{0.1in}

\item \textbf{Delivery Times.} \quad
An online retailer claims the mean delivery time for orders is 3 days. A consumer group
believes the true mean is \emph{different} from 3 days. They randomly audit 55 recent orders.

\begin{enumerate}[label=\alph*.]
  \item Method: \blank{5.5cm} \quad $df = $ \blank{1.5cm}
  \item $H_0$: \blank{4.5cm} \quad $H_a$: \blank{4.5cm}
  \item Conditions:
        \begin{itemize}
          \item Random: \writeline
          \item Normal: \writeline
        \end{itemize}
\end{enumerate}
\vspace{0.1in}

\item \textbf{Matched Pairs: Reaction Time.} \quad
A researcher tests whether a new app reduces reaction time. Ten participants each complete
a reaction-time task \emph{before} and \emph{after} using the app for two weeks. The
differences (before $-$ after) show a roughly symmetric distribution with no outliers.

\begin{enumerate}[label=\alph*.]
  \item Define the order of subtraction and the parameter $\mu_d$:

        $d_i = $ \blank{3cm} $-$ \blank{3cm} \quad; \quad $\mu_d = $ \blank{6cm}

  \item Method: \blank{5cm} \quad $df = $ \blank{1.5cm}
  \item $H_0$: \blank{4cm} \quad $H_a$: \blank{4cm}
  \item Conditions:
        \begin{itemize}
          \item Random/Independence: \writeline
          \item Normal ($n < 30$): \writeline
        \end{itemize}
\end{enumerate}
\vspace{0.1in}

\item \textbf{AP-Style.} \quad
\textit{A researcher collects a random sample of 12 observations to test $H_0: \mu = 50$ vs.\
$H_a: \mu \ne 50$. She verifies that the dotplot of the data shows moderate right skewness and
no outliers. Which of the following best describes the Normal condition?}

\begin{enumerate}[label=(\Alph*)]
  \item The condition is met because $n = 12$ is large enough for the CLT.
  \item The condition is not met; moderate right skewness with $n < 30$ violates the condition.
  \item The condition is met; moderate skewness with no outliers is acceptable for $n < 30$.
  \item The condition cannot be evaluated without knowing the population standard deviation.
\end{enumerate}
\vspace{0.06in}
Answer: \blank{1.5cm} \quad Brief justification: \writeline

\end{enumerate}

\begin{extensionbox}
A statistician says: ``When $n$ is small and the data are not Normal, I can use a sign test or
Wilcoxon signed-rank test instead of a $t$-test.'' Research what a \textbf{nonparametric test}
is and explain in 2--3 sentences why these alternatives exist and what trade-off they involve.
\writelines{3}
\end{extensionbox}

\end{document}
